% Problem formulation (condensed for the two-column paper, step 2). Matches gat_imoea_dqn/evaluator.py,
% intervals.py and dataset_tools/loader.py with EvaluatorConfig(energy_scope="device", external_io="device",
% makespan="sum_apps").
\section{Problem formulation}
\label{sec:model}

\subsection{Interval numbers}
\label{sec:intervals}
An uncertain quantity is an interval $\tilde a = [a^L, a^U]$ with midpoint $m(\tilde a)$ and width $w(\tilde a) =
a^U - a^L$. Every operation of the model (sum, maximum, and a positive quantity divided by a positive interval) is
monotone in each bound, so it is applied to the lower and to the upper inputs, e.g.\ $\tilde C / \tilde f = [C^L / f^U,
C^U / f^L]$. Intervals are compared by the possibility degree \citep{jin2026imoeadqn}
\begin{equation}
P(\tilde a \le \tilde b) = \min\Big\{1, \max\Big\{0, \frac{b^U - a^L}{w(\tilde a) + w(\tilde b)}\Big\}\Big\}.
\label{eq:possibility}
\end{equation}
Solution $x$ \emph{interval-dominates} $y$ ($x \prec_{IN} y$) when $P(\tilde F_m(x) \le \tilde F_m(y)) \ge 0.5$ for every
objective $m$, with strict inequality for at least one. The uncertainty of a solution is measured by $U_1 = \sum_m
w(\tilde F_m)$ and $U_2 = \prod_m w(\tilde F_m)$ on normalised objectives.

\subsection{System and application model}
The system (Fig.~\ref{fig:system_model}) follows \citet{jin2026imoeadqn}: $D$ mobile devices, $n$ edge servers and a
cloud. Device $d$ has CPU frequency $\tilde f_d$, energy per cycle $\tilde e_d$, transmit power $p_d$ and idle power
$p^{idle}_d$, and reaches the edge servers within its range at the wireless rate $\tilde R_{d,u} = \tilde B \log_2(1 + p_d
g / v_0)$. Edge server $u$ shares its frequency $\tilde f_u$ equally among the tasks placed on it, and the cloud gives
every task $\tilde f_C$. Edge servers are linked to each other and to the cloud at rates $\tilde W_E$ and $\tilde W_C$; a
device reaches the cloud through its nearest edge server.

Each device submits one application, a DAG of tasks; together they form $\mathcal{G} = (V, E)$ with $K = |V|$. Task $k$
has workload $\tilde C_k$ cycles, arc $(p, k)$ carries $\tilde D_{pk}$ bits, and a task may read input from and return
output to its device. A task starts only after all its predecessors have finished and their data have arrived. With one
task per application the model reduces to that of \citet{jin2026imoeadqn}.

A solution is a pair $(\mathbf{x}, \mathbf{o})$. The allocation string $\mathbf{x} \in \{0, \dots, n+1\}^K$ places each
task on its device ($0$), an edge server ($1..n$) or the cloud ($n+1$); the priority string $\mathbf{o}$ is a topological
order of $\mathcal{G}$ that fixes the dispatch order.

\subsection{Interval critical-path timing}
Task $k$ runs at $\tilde f_k = \tilde f_d$, $\tilde f_u / n_u$ or $\tilde f_C$ for $n_u$ tasks on server $u$, so it
computes for $\tilde t_k = \tilde C_k / \tilde f_k$. Visiting the tasks in the order $\mathbf{o}$,
\begin{align}
\widetilde{AT}_k &= \max\Big\{\tilde\tau^{in}_k, \max_{(p, k) \in E} \Big(\widetilde{FT}_p + \frac{\tilde D_{pk}}{\tilde R(x_p, x_k)}\Big)\Big\}, \label{eq:at}\\
\widetilde{FT}_k &= \max\{\widetilde{AT}_k, \widetilde{av}(x_k)\} + \tilde t_k, \label{eq:ft}
\end{align}
where $\tilde\tau^{in}_k$ is the upload time of $k$'s input if $k$ is offloaded, transfers between co-located tasks
take no time, and $\widetilde{av}(x_k)$ is the finish time of the previous task on the same device, since a device runs
one task at a time. Maxima are taken boundwise; the chain that sets the latest finish time is the \emph{interval
critical path}.

\subsection{Objectives and problem}
The objectives are the total makespan, the devices' energy and the monetary cost:
\begin{align}
\tilde F_1 &= \sum_{a} \max_{k \in V_a} \big(\widetilde{FT}_k + \tilde\tau^{out}_k\big), \label{eq:f1}\\
\tilde F_2 &= \sum_{x_k = 0} \tilde C_k \tilde e_{d(k)} \nonumber\\
  &\quad + \sum_{x_k \ne 0} \Big(p_{d(k)} \frac{\tilde D^{in}_k}{\tilde R_{d(k), x_k}} + p^{idle}_{d(k)} \tilde t_k\Big) \nonumber\\
  &\quad + \sum_{\substack{(p, k) \in E:\\ x_p = 0,\, x_k \ne 0}} p_{d(p)} \frac{\tilde D_{pk}}{\tilde R_{d(p), x_k}}, \label{eq:f2}\\
\tilde F_3 &= \eta_E \sum_{1 \le x_k \le n} \tilde f_k \tilde t_k + \eta_C \sum_{x_k = n+1} \tilde t_k, \label{eq:f3}
\end{align}
where $V_a$ are the tasks of application $a$, $\tilde\tau^{out}_k$ the time to return $k$'s output and $d(k)$ its
device. $\tilde F_2$ counts local computation, transmission to offloaded tasks (for the cloud, up to the nearest edge
server, idling while the data cross $\tilde W_C$) and idle waiting, i.e.\ the devices'
energy only, as in \citet{jin2026imoeadqn} (all-tier energy is a sensitivity case in Section~\ref{sec:sensitivity});
$\eta_E$ and $\eta_C$ are the edge and cloud prices. The problem is
\begin{equation}
\min_{(\mathbf{x}, \mathbf{o})} \big(\tilde F_1, \tilde F_2, \tilde F_3\big)\ \text{s.t.}\ x_k \in \mathcal{A}_k\, \forall k,\
\mathbf{o}\ \text{topol.},
\label{eq:problem}
\end{equation}
where $\mathcal{A}_k$ are the nodes reachable from $k$'s device; we seek the solutions that no other solution
interval-dominates. Scheduling a DAG on heterogeneous processors is NP-hard even with crisp data
\citep{topcuoglu2002heft}, which motivates an evolutionary approach.
